\documentclass[12pt,a4paper]{report}
\usepackage[utf8]{inputenc}
\usepackage[T1]{fontenc}
\usepackage[french]{babel}
\usepackage{palatino}
\usepackage{geometry}
\geometry{
    a4paper,
    top=2.5cm,
    bottom=2.5cm,
    left=3cm,
    right=2.5cm
}
\usepackage{graphicx}
\usepackage[colorlinks=true, linkcolor=black, citecolor=red]{hyperref}
\usepackage{fancyhdr}
\usepackage{enumitem}

% En-têtes et pieds de page
\pagestyle{fancy}
\fancyhf{}
\fancyhead[L]{\nouppercase{\leftmark}}
\fancyhead[R]{\thepage}
\renewcommand{\headrulewidth}{0.4pt}

% Pour la table des matières
\usepackage{tocloft}
\setlength{\cftchapnumwidth}{2em}
\setlength{\cftsecnumwidth}{2.5em}

\begin{document}

% ========== PAGE DE TITRE ==========
\begin{titlepage}
\centering
\vspace*{2cm}
{\LARGE \textbf{Plan du Mémoire de Fin de Formation}} \\[0.8cm]
{\Large \textbf{Plateforme d’orientation professionnelle — orientation-bj}} \\[2cm]
{\large \textbf{Présenté par :}} \\[0.3cm]
{\large D'ALMEIDA Gédéon} \\
{\large METINHOUE Paola} \\[2cm]
{\large \today}
\end{titlepage}

% ========== TABLE DES MATIÈRES ==========
\tableofcontents
\newpage

% ========== INTRODUCTION (sans numéro mais dans la TOC) ==========
\chapter*{Introduction}
\addcontentsline{toc}{chapter}{Introduction}

\section*{1. Contexte et justification}
\addcontentsline{toc}{section}{1. Contexte et justification}
\subsection*{Contexte}
L’orientation scolaire et professionnelle est souvent freinée par un manque d’outils accessibles, d’informations fiables et d’accompagnement continu. Les apprenants et jeunes actifs peinent à relier leurs intérêts aux possibilités réelles de formation et de métiers.

\subsection*{Justification}
Il devient nécessaire de proposer un écosystème unifié qui aide à l’auto‑connaissance, centralise les contenus métiers/formations et facilite une décision d’orientation éclairée.

\section*{2. Problématique}
\addcontentsline{toc}{section}{2. Problématique}
Comment concevoir une plateforme d’orientation professionnelle unifiée, inclusive et fiable, capable d’aider les apprenants et jeunes actifs à mieux se connaître, à explorer les métiers et à prendre des décisions d’orientation éclairées ?

\section*{3. Objectifs du projet de fin de formation}
\addcontentsline{toc}{section}{3. Objectifs du projet de fin de formation}
\subsection*{Objectif général}
Concevoir et présenter l’écosystème d’orientation professionnelle « orientation-bj », fondé sur le modèle RIASEC, afin d’en démontrer la pertinence et l’utilité pour la prise de décision d’orientation.

\subsection*{Objectifs spécifiques}
\begin{itemize}
    \item Concevoir un parcours d’auto‑évaluation basé sur RIASEC et interprétation pédagogique des résultats.
    \item Proposer un système de recommandations métiers et formations adapté au profil.
    \item Mettre à disposition des contenus fiables et contextualisés pour l’orientation.
\end{itemize}

\section*{4. Plan du document}
\addcontentsline{toc}{section}{4. Plan du document}
Annonce textuelle et fluide de l’articulation des chapitres.

\newpage

% ============================================================
% CHAPITRE 1
% ============================================================
\chapter{Déroulement du stage}

\section{Introduction}
Cadre du stage, durée et objectifs personnels en tant que futur concepteur de solutions numériques.

\section{Présentation de la structure}
Fiche d’identité, historique et positionnement de la structure d’accueil.

\section{Fonctionnement de la structure}
Organisation interne et interactions entre les équipes impliquées.

\section{Travaux effectués}
Missions réalisées autour du projet d’orientation RIASEC, organisation du travail et méthodologies adoptées.

\section{Apports du stage sur le plan professionnel}
Compétences acquises, apprentissages et apports méthodologiques.

\section{Difficultés rencontrées et suggestions}
\subsection*{Difficultés}
Contraintes et défis rencontrés durant le stage.
\subsection*{Suggestions}
Améliorations techniques ou organisationnelles proposées.

\section{Conclusion}
Bilan synthétique de l’immersion professionnelle.

% ============================================================
% CHAPITRE 2
% ============================================================
\chapter{Analyse et Conception}

\section{Introduction}
Présentation de la démarche d’ingénierie pour traduire les besoins d’orientation en spécifications.

\section{État de l’art}
\subsection*{Analyse}
Panorama des dispositifs d’orientation, tests psychométriques et plateformes numériques existantes.
\subsection*{Critique}
Limites observées : faible contextualisation, accessibilité inégale, accompagnement insuffisant.

\section{Présentation du projet}
Présentation globale d’« orientation-bj » : un écosystème d’orientation fondé sur RIASEC, destiné aux apprenants et jeunes actifs.

\section{Présentation des fonctionnalités}
Parcours d’évaluation, restitution des résultats, recommandations métiers/formations, et contenus d’accompagnement.

\section{Analyse et modélisation}
Modélisation UML ciblée :
\begin{itemize}
    \item Diagramme de cas d’utilisation pour clarifier les rôles.
    \item Diagramme de classes pour structurer les données principales.
    \item Diagrammes de séquence pour illustrer les interactions clés du parcours.
\end{itemize}

\section{Matériel}
\subsection*{Logiciel}
Environnement applicatif (front, back, base de données), outils d’analyse et de suivi.
\subsection*{Physique}
Postes de test, équipements mobiles et infrastructure réseau de démonstration.

\section{Conclusion}
Validation des besoins comme base de la conception.

% ============================================================
% CHAPITRE 3
% ============================================================
\chapter{Résultats et Discussion}

\section{Introduction}
Présentation de la double démarche : exposer la réalisation et discuter la réponse aux objectifs.

\section{Résultats}
\begin{itemize}
    \item \textbf{Implémentation (captures d’écran légendées)} : parcours d’évaluation, résultats RIASEC et recommandations.
    \item \textbf{Preuve de cohérence} : exemples de correspondances profil--métier/formation.
    \item \textbf{Fiabilité du parcours} : clarté des résultats et logique d’accompagnement.
\end{itemize}

\section{Discussions}
\begin{itemize}
    \item \textbf{Apports du projet} : meilleure lisibilité des choix d’orientation.
    \item \textbf{Limites et arbitrages} : contraintes de données, disponibilité des contenus, accompagnement humain.
\end{itemize}

\section{Conclusion}
Constat de la conformité du produit final par rapport aux objectifs fixés.

\section{Gouvernance des contenus}
Sources, validation, mise à jour et traçabilité des informations.

\section{Sécurité et confidentialité}
Consentement, droits d’accès et protection des données personnelles.

\section{Conclusion}
Cohérence de la conception au regard des exigences.

% ============================================================
% CHAPITRE 4
% ============================================================
% ============================================================
% CONCLUSION GÉNÉRALE ET PERSPECTIVES
% ============================================================
\chapter*{Conclusion générale et Perspectives}
\addcontentsline{toc}{chapter}{Conclusion générale et Perspectives}

\section*{1. Conclusion}
Synthèse des contributions du travail et bilan du projet d’orientation RIASEC.

\section*{2. Perspectives}
\begin{itemize}
    \item Extension du contenu métiers et formations avec des partenaires éducatifs.
    \item Amélioration des parcours d’accompagnement et de suivi.
    \item Ouverture à des usages institutionnels (établissements, conseillers, employeurs).
\end{itemize}

\end{document}
